% Teil II: der durchgerechnete Weg.

\subsection*{Gegeben und gesucht}
(a) Zugverspätung mit $p=\tfrac15$ je Fahrt; gesucht:
$\mathrm P(\text{mindestens 4 von 5 verspätet})$.
(b) Aquarium mit $9$ Fischen, je $p=\tfrac1{10}$ pro Nacht;
gesucht: $\mathrm P(\text{mindestens einer springt hinaus})$.

\subsection*{a) Mindestens 4 von 5 verspätet}
$k=4$ oder $k=5$:
\begin{align}
  \mathrm P &= \binom54\Bigl(\frac15\Bigr)^4\frac45
     +\binom55\Bigl(\frac15\Bigr)^5 \\
  &=5\cdot\frac{1}{625}\cdot\frac45+\frac1{3125}
   =\frac{20}{3125}+\frac1{3125}=\frac{21}{3125}.
\end{align}
\begin{ergebnis}
$\mathrm P=\dfrac{21}{3125}\approx0.0067$.
\end{ergebnis}

\subsection*{b) Mindestens einer von 9 springt hinaus}
Gegenereignis ``keiner springt hinaus'':
\begin{equation}
  \mathrm P=1-\Bigl(1-\frac1{10}\Bigr)^9
  =1-\Bigl(\frac9{10}\Bigr)^9 .
\end{equation}
\begin{ergebnis}
$\mathrm P=1-\bigl(\tfrac9{10}\bigr)^9
=1-\tfrac{387420489}{10^9}\approx0.6126$.
\end{ergebnis}

\begin{kontrolle}
a) $p$ klein, $k\ge4$: Ergebnis $~0.0067$ ist klein aber positiv;
b) $\bigl(\tfrac9{10}\bigr)^9\approx0.387$, also $P\approx0.613<1$ -- plausibel.
\end{kontrolle}

\begin{merksatz}
Bei ``mindestens $c$'' summiert man die Binomialwahrscheinlichkeiten ab
$c$; bei ``mindestens einer'' ist das Gegenereignis ``keiner'' meist
schneller.
\end{merksatz}
